\section{Implication for a Parameterized Matroid Problem}
\label{sec:additional}


In this Section we describe  an additional implication of \Cref{thm:emi} for a parameterized decision problem considered by Fomin et al.\ \cite{FGKSS23}. In \cite{FGKSS23} the authors present a ``pseudo-polynomial parameterized algorithm'' for a matroid problem in which the matroid is representable over a finite field. We use \Cref{thm:emi} to show the result cannot be extended to arbitrary matroids in the oracle model while retaining a similar running time. 


Given a graph $G=(V,F)$ and a path $Q$ in $G$ we use $V(Q)$ to denote the set of vertices on the path $Q$.   Consider the following problem.
%Consider for example {\em generalized weighted colored paths} problem considered by Fomin et al.\ \cite{FGKSS23}.

\problem{  {\sc Generalized Weighted Colored Paths} (GWCP)}{
	{\bf Instance }  & An undirected  graph $G=(V,F)$, two sets $A,B\subseteq V$, a representable matroid $(V, \cI)$ over a finite field of order $q$ (that is, the elements of the matroid are the vertices of the graph), and $p,k,T\in \mathbb{N}$.\\
	{\bf A solution} & $p$ vertex disjoint paths $Q_1,\ldots, Q_p$  starting from a vertex in  $A$ and terminating in a vertex in  $B$, such that there is $X\subseteq V(Q_1)\cup \ldots \cup V(Q_p)$  which satisfies $X\in \cI$, $w(X)=T$ and $|X|=k$. \\
	{\bf Objective}& Decide if there is a solution.
}

In \cite{FGKSS23} the authors proposed an algorithm for GWCP which runs in time $2^{p+O(k^2 \log(q+k ))} \cdot n^{O(1)}\cdot T$, where $n$ is the number of vertices of the input graph (Theorem 1.5 in \cite{FGKSS23}, the problem was not given a name in \cite{FGKSS23}).  
Similar to the algorithm of \cite{CGM92} for representable EMI, the algorithm of~\cite{FGKSS23} is algebraic and relies on the DeMillo-Lipton-Schwartz-Zippel lemma for polynomial identity testing~\cite{Sc80,DL78,Zi79}.
 Consider the following variant of GWCP, in which the matroid is arbitrary and given via a membership oracle. 

\problem{  {\sc Oracle Generalized Weighted Colored Paths} (Oracle-GWCP)}{
	{\bf Instance }  & $(G=(V,F),A,B,\cI,p,k,T)$ where  $G=(V,F)$ is an undirected graph, $A,B\subseteq V$, $(V, \cI)$ is a matroid, and $p,k,T\in \mathbb{N}$.\\
	{\bf A solution} & $p$ vertex disjoint paths $Q_1,\ldots, Q_p$ starting from a vertex in  $A$ and terminating in a vertex in  $B$, such that that there is $X\subseteq V(Q_1)\cup \ldots \cup V(Q_p)$  which satisfies $X\in \cI$, $w(X)=T$ and $|X|=k$. \\
	{\bf Objective}& Decide if there is a solution.\\
	{\bf Model}& The independent sets of the matroid are accessible only via a membership oracle. 
}


We can use \Cref{thm:emi} to show a lower bound on the running time of algorithms for  Oracle-GWCP. 
\begin{lemma}
There is no randomized algorithm for \textnormal{Oracle-GWCP}  which runs in time $g(k,p)\cdot (n\cdot T)^{O(1)}$ for some function $g(k,p)$.
\end{lemma}
\begin{proof}
Assume towards contradiction that there is a randomized algorithm  $\cA$ for Oracle-GWCP   which runs in time $g(k,p)\cdot (n\cdot T)^{O(1)}$ for some function $g(k,p)$.  Consider the following algorithm for EMI, to which we refer as $\cB$, whose input is an EMI instance $I=(E,\cI,w,k,T)$.
\begin{itemize}
	\item If $|E|<2$ then solve the instance checking for every $S\subseteq E$ if it is a solution.
	\item   Let $E=\{e_1,\ldots, e_n\}$;  define $A=\{e_1\}$, $B=\{e_n\}$, $F=\{ (e_i,e_{i+1})~|~1\leq i<n\}$, and $p=1$. 
	\item  Use $\cA$ to solve the Oracle-GWCP instance $J=(G=(E,F),A,B,\cI, p,k,T)$. If $\cA$ determines  that $J$ is a ``yes''-instance then return that $I$ is a yes-instance. Otherwise, return that $J$ is a ``no''-instance. 
\end{itemize}
Observe that the graph $G=(E,F)$ is a path from $e_1$ to $e_n$ which traverses all the elements (vertices) in $E$. 
Consider the execution of the algorithm given an instance $I=(E,\cI,w,k,T)$. Clearly, the algorithm returns a correct solution in case $|E| <2$. Hence, we may assume w.l.o.g. that $|E|\geq 2$.  Let~$J$ be the Oracle-GWCP instance constructed by the algorithm.

\begin{claim}
	\label{claim:gwcp}
	The \textnormal{EMI} instance $I$  is a ``yes''-instance iff  the \textnormal{Oracle-GWCP} instance $J$ is a ``yes''-instance.
\end{claim}
\begin{claimproof}
	If $I$ is a ``yes''-instance, let $X$ be its solution. Also, let $Q_1$ be the single path in $G$ from $e_1$ to $e_n$. It holds that  $X\subseteq V(Q_1)=E$. Since $X$ is a solution for $I$, we have $X\in \cI$, $|X|=k$ and $w(X)=T$. That is, $Q_1$ is a solution for $J$, and therefore $J$ is a yes instance.
	
	If $J$ is a ``yes''-instance then there is a path $Q_1$ from $e_1$ to $e_n$ and $X\subseteq V(Q_1)$ such that $X\in \cI$, $|X|=k$ and $w(X)=T$. Thus, $X$ is a solution for $I$, and therefore $I$ is a ``yes''-instance.
\end{claimproof}

By \Cref{claim:gwcp}, and since $\cA$ is a randomized algorithm for Oracle-GWCP, it holds that $\cB$ is a randomized algorithm for EMI. The running time of $\cB$ is $g(k,1)\cdot (n\cdot T)^{O(1)}$, in contradiction to \Cref{thm:emi}. Thus, there is no randomized algorithm for Oracle-GWCP which runs in time $g(k,p)\cdot (n\cdot T)^{O(1)}$. 
	
\end{proof}

%
%It can be easily observed that an algorithm for Oracle-GWCP with running time of the form $g(k,p) \cdot \left(n\cdot T\right)^{O(1)}$  for an arbitrary function $g$ can be used to decide EMI in time $f(k) \cdot \left(|I|\cdot T\right)^{O(1)}$. Given an EMI instance $(E,\cI,w,k,T)$  we can generate a graph $G=(E,F)$ where $F$ is a path between two arbitrary elements $a,b\in E$. Also, define $A=\{a\}$ and $B=\{b\}$. It can be easily verified that this Oracle-GWCP has a solution if and only if the EMI instance is a yes instance. Thus, an Oracle-GWCP algorithm in time  $g(k,p) \cdot \left(n\cdot T\right)^{O(1)}$ implies an EMI algorithm in time $g(k,1)\cdot (|I| \cdot T)^{O(1)}$,  in contradiction to  \Cref{thm:emi}.